\documentclass[aspectratio=169, xcolor=dvipsnames]{beamer}

% --- Theme Setup ---
\usetheme{Madrid}
\usecolortheme{default}

% --- Packages ---
\usepackage{graphicx}
\usepackage{xcolor}
\usepackage{booktabs}
\usepackage{hyperref}
\usepackage{adjustbox}
\usepackage{amssymb}
\usepackage{amsmath}
\usepackage{listings}
\usepackage{caption}
\usepackage{tikz}

\lstset{
  basicstyle=\ttfamily\small,
  keywordstyle=\color{blue}\bfseries,
  stringstyle=\color{red},
  showstringspaces=false,
  breaklines=true
}

% --- Title Information ---
\title{Module 25}
\subtitle{Relational Database Design/5}
\author{Milav Dabgar}
\institute{www.milav.in}
\date{\today}

\begin{document}

% Title Slide
\begin{frame}
    \titlepage
\end{frame}

\begin{frame}{Module Recap}
    \begin{itemize}[<+->]
        \item Studied Algorithms for Properties of Functional Dependencies
    \end{itemize}
\end{frame}

\begin{frame}{Module Objectives}
    \begin{block}{Goals}
        \begin{itemize}[<+->]
            \item To Understand the Characterizations for Lossless Join Decomposition
            \item To Understand the Characterizations for Dependency Preservation
        \end{itemize}
    \end{block}
\end{frame}

\begin{frame}{Module Outline}
    \begin{itemize}[<+->]
        \item Lossless Join Decomposition
        \item Dependency Preservation
    \end{itemize}
\end{frame}

\section{Lossless Join Decomposition}
\begin{frame}
  \vfill
  \centering
  \begin{beamercolorbox}[sep=8pt,center,shadow=true,rounded=true]{title}
    \usebeamerfont{title}Lossless Join Decomposition\par%
  \end{beamercolorbox}
  \vfill
\end{frame}

\begin{frame}{Lossless Join Decomposition}
    \begin{itemize}[<+->]
        \item For the case of $R = (R_1, R_2)$, we require that for all possible relations $r$ on schema $R$
        \[ r = \Pi_{R_1}(r) \bowtie \Pi_{R_2}(r) \]
        \item A decomposition of $R$ into $R_1$ and $R_2$ is lossless join if at least one of the following dependencies is in $F^+$:
        \begin{itemize}
            \item[$\triangleright$] $R_1 \cap R_2 \rightarrow R_1$
            \item[$\triangleright$] $R_1 \cap R_2 \rightarrow R_2$
        \end{itemize}
        \item The above functional dependencies are a sufficient condition for lossless join decomposition; the dependencies are a necessary condition only if all constraints are functional dependencies
    \end{itemize}
\end{frame}

\begin{frame}{Lossless Join Decomposition (continued)}
    To Identify whether a decomposition is lossy or lossless, it must satisfy the following conditions:
    \begin{itemize}[<+->]
        \item $R_1 \cup R_2 = R$
        \item $R_1 \cap R_2 \neq \phi$ and
        \item $R_1 \cap R_2 \rightarrow R_1$ or $R_1 \cap R_2 \rightarrow R_2$
    \end{itemize}
\end{frame}

\begin{frame}{Lossless Join Decomposition (2): Example}
    \begin{alertblock}{Lossy Decomposition Example}
        \begin{itemize}[<+->]
            \item Consider Supplier Parts schema: \texttt{Supplier\_Parts(S\#, Sname, City, P\#, Qty)}
            \item Having dependencies: \texttt{S\# $\rightarrow$ Sname, S\# $\rightarrow$ City, (S\#, P\#) $\rightarrow$ Qty}
            \item Decompose as: \texttt{Supplier(S\#, Sname, City, Qty) : Parts(P\#, Qty)}
            \item Take Natural Join to reconstruct: \texttt{Supplier $\bowtie$ Parts}
            \item We get extra tuples! Join is \textbf{Lossy}!
            \item Common attribute \texttt{Qty} is not a superkey in Supplier or in Parts
            \item Does not preserve \texttt{(S\#, P\#) $\rightarrow$ Qty}
        \end{itemize}
    \end{alertblock}
\end{frame}

\begin{frame}{Lossless Join Decomposition (3): Example}
    \begin{exampleblock}{Lossless Decomposition Example}
        \begin{itemize}[<+->]
            \item Consider Supplier Parts schema: \texttt{Supplier\_Parts(S\#, Sname, City, P\#, Qty)}
            \item Having dependencies: \texttt{S\# $\rightarrow$ Sname, S\# $\rightarrow$ City, (S\#, P\#) $\rightarrow$ Qty}
            \item Decompose as: \texttt{Supplier(S\#, Sname, City) : Parts(S\#, P\#, Qty)}
            \item Take Natural Join to reconstruct: \texttt{Supplier $\bowtie$ Parts}
            \item We get back the original relation. Join is \textbf{Lossless}.
            \item Common attribute \texttt{S\#} is a superkey in Supplier
            \item Preserves all dependencies
        \end{itemize}
    \end{exampleblock}
\end{frame}

\begin{frame}{Lossless Join Decomposition (4): Example}
    \begin{itemize}[<+->]
        \item $R = (A, B, C)$
        \item $F = \{A \rightarrow B, B \rightarrow C\}$
        \item Can be decomposed in two different ways
        \item $R_1 = (A, B), R_2 = (B, C)$
        \begin{itemize}
            \item[$\triangleright$] Lossless-join decomposition: $R_1 \cap R_2 = \{B\}$ and $B \rightarrow BC$
            \item[$\triangleright$] Dependency preserving
        \end{itemize}
        \item $R_1 = (A, B), R_2 = (A, C)$
        \begin{itemize}
            \item[$\triangleright$] Lossless-join decomposition: $R_1 \cap R_2 = \{A\}$ and $A \rightarrow AB$
            \item[$\triangleright$] Not dependency preserving (cannot check $B \rightarrow C$ without computing $R_1 \bowtie R_2$)
        \end{itemize}
    \end{itemize}
\end{frame}

\begin{frame}{Practice Problems on Lossless Join}
    \begin{itemize}[<+->]
        \item Check if the decomposition of $R$ into $D$ is lossless:
        \item a) $R(ABC)$: $F = \{A \rightarrow B, A \rightarrow C\}$. $D = R_1(AB), R_2(BC)$
        \item b) $R(ABCDEF)$: $F = \{A \rightarrow B, B \rightarrow C, C \rightarrow D, E \rightarrow F\}$. $D = R_1(AB), R_2(BCD), R_3(DEF)$
        \item c) $R(ABCDEF)$: $F = \{A \rightarrow B, C \rightarrow DE, AC \rightarrow F\}$. $D = R_1(BE), R_2(ACDEF)$
        \item d) $R(ABCDEG)$: $F = \{AB \rightarrow C, AC \rightarrow B, AD \rightarrow E, B \rightarrow D, BC \rightarrow A, E \rightarrow G\}$
        \begin{itemize}
            \item[$\triangleright$] i) $D_1 = R_1(AB), R_2(BC), R_3(ABDE), R_4(EG)$
            \item[$\triangleright$] ii) $D_2 = R_1(ABC), R_2(ACDE), R_3(ADG)$
        \end{itemize}
        \item e) $R(ABCDEFGHIJ)$: $F = \{AB \rightarrow C, B \rightarrow F, D \rightarrow IJ, A \rightarrow DE, F \rightarrow GH\}$
        \begin{itemize}
            \item[$\triangleright$] i) $D_1 = R_1(ABC), R_2(ADE), R_3(BF), R_4(FGH), R_5(DIJ)$
            \item[$\triangleright$] ii) $D_2 = R_1(ABCDE), R_2(BFGH), R_3(DIJ)$
            \item[$\triangleright$] iii) $D_3 = R_1(ABCD), R_2(DE), R_3(BF), R_4(FGH), R_5(DIJ)$
        \end{itemize}
    \end{itemize}
\end{frame}

\section{Dependency Preservation}
\begin{frame}
  \vfill
  \centering
  \begin{beamercolorbox}[sep=8pt,center,shadow=true,rounded=true]{title}
    \usebeamerfont{title}Dependency Preservation\par%
  \end{beamercolorbox}
  \vfill
\end{frame}

\begin{frame}{Dependency Preservation}
    \begin{block}{Definition}
        \begin{itemize}[<+->]
            \item Let $F_i$ be the set of dependencies $F^+$ that include only attributes in $R_i$
            \item A decomposition is \textbf{dependency preserving}, if
            \[ (F_1 \cup F_2 \cup \dots \cup F_n)^+ = F^+ \]
            \item If it is not, then checking updates for violation of functional dependencies may require computing joins, which is expensive
        \end{itemize}
    \end{block}
\end{frame}

\begin{frame}{Dependency Preservation (continued)}
    Let $R$ be the original relational schema having FD set $F$. Let $R_1$ and $R_2$ having FD set $F_1$ and $F_2$ respectively, are the decomposed sub-relations of $R$. The decomposition of $R$ is said to be preserving if
    \begin{itemize}[<+->]
        \item $F_1 \cup F_2 \equiv F$ \{ Decomposition Preserving Dependency \}
        \item If $F_1 \cup F_2 \subset F$ \{ Decomposition NOT Preserving Dependency \} and
        \item $F_1 \cup F_2 \supset F$ \{ this is not possible \}
    \end{itemize}
\end{frame}

\begin{frame}[fragile]{Dependency Preservation (2): Testing}
    \begin{itemize}[<+->]
        \item To check if a dependency $\alpha \rightarrow \beta$ is preserved in a decomposition of $R$ into $D = \{R_1, R_2, \dots, R_n\}$ we apply the following test (with attribute closure done with respect to $F$)
        \item The restriction of $F^+$ to $R_i$ is the set of all functional dependencies in $F^+$ that include only attributes of $R_i$.
    \end{itemize}
    \vspace{0.5em}
    \begin{verbatim}
compute F+;
for each schema Ri in D do begin
    Fi = the restriction of F+ to Ri;
end
F' = phi
for each restriction Fi do begin
    F' = F' union Fi
end
compute F'+;
if (F'+ == F+) then return (true) else return (false);
    \end{verbatim}
    \begin{itemize}[<+(2)->]
        \item The procedure for checking dependency preservation takes exponential time to compute $F^+$ and $(F_1 \cup F_2 \cup \dots \cup F_n)^+$
    \end{itemize}
\end{frame}

\begin{frame}{Dependency Preservation (3): Example}
    \begin{itemize}[<+->]
        \item $R(A, B, C, D, E, F)$
        \item $F = \{A \rightarrow BCD, A \rightarrow EF, BC \rightarrow AD, BC \rightarrow E, BC \rightarrow F, B \rightarrow F, D \rightarrow E\}$
        \item Decomposition:
        \item $R_1(A, B, C, D) \qquad R_2(B, F) \qquad R_3(D, E)$
        \item $A \rightarrow BCD, BC \rightarrow AD$ are preserved on table $R_1$
        \item $D \rightarrow E$ is preserved on table $R_3$
        \item $B \rightarrow F$ is preserved on table $R_2$
        \item We have to check whether the remaining FDs: $A \rightarrow E, A \rightarrow F, BC \rightarrow E, BC \rightarrow F$ are preserved or not.
    \end{itemize}
\end{frame}

\begin{frame}{Dependency Preservation (3): Example - Checking}
    \begin{itemize}[<+->]
        \item $F' = F_1 \cup F_2 \cup F_3$
        \item Checking for: $A \rightarrow E, A \rightarrow F$ in $F'^+$
        \begin{itemize}
            \item[$\triangleright$] $A \rightarrow D$ (from $R_1$), $D \rightarrow E$ (from $R_3$) : $A \rightarrow E$ (By Transitivity)
            \item[$\triangleright$] $A \rightarrow B$ (from $R_1$), $B \rightarrow F$ (from $R_2$) : $A \rightarrow F$ (By Transitivity)
        \end{itemize}
        \item Checking for: $BC \rightarrow E, BC \rightarrow F$ in $F'^+$
        \begin{itemize}
            \item[$\triangleright$] $BC \rightarrow D$ (from $R_1$), $D \rightarrow E$ (from $R_3$) : $BC \rightarrow E$ (By Transitivity)
            \item[$\triangleright$] $BC \rightarrow F$ (By Augmentation on $B \rightarrow F$ from $R_2$)
        \end{itemize}
    \end{itemize}
    \vspace{1em}
    \uncover<5->{\textbf{Hence all dependencies are preserved.}}
\end{frame}

\begin{frame}{Dependency Preservation (4): Example}
    \begin{itemize}[<+->]
        \item $R(A, B, C, D)$
        \item $F = \{A \rightarrow B, B \rightarrow C, C \rightarrow D, D \rightarrow A\}$
        \item Decomposition:
        \item $R_1(A, B) \qquad R_2(B, C) \qquad R_3(C, D)$
        \item $A \rightarrow B$ is preserved on table $R_1$
        \item $B \rightarrow C$ is preserved on table $R_2$
        \item $C \rightarrow D$ is preserved on table $R_3$
        \item We have to check whether the one remaining FD: $D \rightarrow A$ is preserved or not.
        \item $F' = F_1 \cup F_2 \cup F_3$.
        \item Checking for: $D \rightarrow A$ in $F'^+$
        \begin{itemize}
            \item[$\triangleright$] $D \rightarrow C$ (from $R_3$), $C \rightarrow B$ (from $R_2$), $B \rightarrow A$ (from $R_1$) : $D \rightarrow A$ (By Transitivity)
        \end{itemize}
    \end{itemize}
    \vspace{1em}
    \uncover<9->{\textbf{Hence all dependencies are preserved.}}
\end{frame}

\begin{frame}[fragile]{Dependency Preservation (5): Testing}
    \begin{itemize}[<+->]
        \item To check if a dependency $\alpha \rightarrow \beta$ is preserved in a decomposition of $R$ into $R_1, R_2, \dots, R_n$ we apply the following test (with attribute closure done with respect to $F$)
    \end{itemize}
    \vspace{0.5em}
    \begin{verbatim}
result = alpha
while (changes to result) do
    for each R_i in the decomposition
        t = (result intersect R_i)^+ intersect R_i
        result = result union t
    \end{verbatim}
    \begin{itemize}[<+(2)->]
        \item If \texttt{result} contains all attributes in $\beta$, then the functional dependency $\alpha \rightarrow \beta$ is preserved.
        \item We apply the test on all dependencies in $F$ to check if a decomposition is dependency preserving
        \item This procedure takes polynomial time, instead of the exponential time required to compute $F^+$ and $(F_1 \cup F_2 \cup \dots \cup F_n)^+$
    \end{itemize}
\end{frame}

\begin{frame}{Dependency Preservation (6): Example}
    \begin{itemize}[<+->]
        \item $R(ABCDEF)$: $F = \{A \rightarrow BCD, A \rightarrow EF, BC \rightarrow AD, BC \rightarrow E, BC \rightarrow F, B \rightarrow F, D \rightarrow E\}$
        \item Decomp = $\{ABCD, BF, DE\}$
        \item Need to check for: $\bigstar$ $A \rightarrow BCD, A \rightarrow EF$, $\bigstar$ $BC \rightarrow AD, BC \rightarrow E, BC \rightarrow F$, $\times$ $B \rightarrow F$, $\times$ $D \rightarrow E$
        \item $(BC)^+ / F_1 = ABCD$. $(ABCD)^+ / F_2 = ABCDF$. $(ABCDF)^+ / F_3 = ABCDEF$. \textbf{Preserves} $BC \rightarrow E, BC \rightarrow F$
        \begin{itemize}
            \item[$\triangleright$] $BC \rightarrow AD (R_1), AD \rightarrow E (R_3)$ implies $BC \rightarrow E$
            \item[$\triangleright$] $B \rightarrow F (R_2)$ implies $BC \rightarrow F$
        \end{itemize}
        \item $(A)^+ / F_1 = ABCD$. $(ABCD)^+ / F_2 = ABCDF$. $(ABCDF)^+ / F_3 = ABCDEF$. \textbf{Preserves} $A \rightarrow EF$
        \begin{itemize}
            \item[$\triangleright$] $A \rightarrow B (R_1), B \rightarrow F (R_2)$ implies $A \rightarrow F$
            \item[$\triangleright$] $A \rightarrow D (R_1), D \rightarrow E (R_3)$ implies $A \rightarrow E$
        \end{itemize}
    \end{itemize}
\end{frame}

\begin{frame}{Dependency Preservation (7): Example}
    \begin{itemize}[<+->]
        \item $R(ABCDEF)$: $F = \{A \rightarrow BCD, A \rightarrow EF, BC \rightarrow AD, BC \rightarrow E, BC \rightarrow F, B \rightarrow F, D \rightarrow E\}$. Decomp = $\{ABCD, BF, DE\}$
        \item Infer reverse FD's:
        \begin{itemize}
            \item[$\triangleright$] $B^+ / F = BF: B \rightarrow A$ cannot be inferred
            \item[$\triangleright$] $C^+ / F = C: C \rightarrow A$ cannot be inferred
            \item[$\triangleright$] $D^+ / F = DE: D \rightarrow A$ and $D \rightarrow BC$ cannot be inferred
            \item[$\triangleright$] $A^+ / F = ABCDEF: A \rightarrow BC$ can be inferred, but it is equal to $A \rightarrow B$ and $A \rightarrow C$
            \item[$\triangleright$] $F^+ / F = F: F \rightarrow B$ cannot be inferred
            \item[$\triangleright$] $E^+ / F = E: E \rightarrow D$ cannot be inferred
        \end{itemize}
        \item $(BC)^+ / F = ABCDEF$. \textbf{Preserves} $BC \rightarrow E, BC \rightarrow F$
        \item $(A)^+ / F = ABCDEF$. \textbf{Preserves} $A \rightarrow EF$
    \end{itemize}
\end{frame}

\begin{frame}{Practice Problems on Dependency Preservation}
    \begin{itemize}[<+->]
        \item Check whether the decomposition of $R$ into $D$ is preserving dependency:
        \item a) $R(ABCD)$: $F = \{A \rightarrow B, B \rightarrow C, C \rightarrow D, D \rightarrow A\}$. $D = \{AB, BC, CD\}$
        \item b) $R(ABCDEF)$: $F = \{AB \rightarrow CD, C \rightarrow D, D \rightarrow E, E \rightarrow F\}$. $D = \{AB, CDE, EF\}$
        \item c) $R(ABCDEG)$: $F = \{AB \rightarrow C, AC \rightarrow B, BC \rightarrow A, AD \rightarrow E, B \rightarrow D, E \rightarrow G\}$. $D = \{ABC, ACDE, ADG\}$
        \item d) $R(ABCD)$: $F = \{A \rightarrow B, B \rightarrow C, C \rightarrow D, D \rightarrow B\}$. $D = \{AB, BC, BD\}$
        \item e) $R(ABCDE)$: $F = \{A \rightarrow BC, CD \rightarrow E, B \rightarrow D, E \rightarrow A\}$. $D = \{ABCE, BD\}$
    \end{itemize}
\end{frame}

\begin{frame}{Module Summary}
    \begin{block}{Summary}
        \begin{itemize}[<+->]
            \item Understood the Characterization for and Determination of Lossless Join
            \item Understood the Characterization for and Determination of Dependency Preservation
        \end{itemize}
    \end{block}
    \vspace{2em}
    \begin{center}
    \small \textit{Slides used in this presentation are borrowed from \url{http://db-book.com/} with kind permission of the authors.} \\
    \small \textit{Edited and new slides are marked with 'PPD'.}
    \end{center}
\end{frame}

\end{document}
